\section{Introduction}
Exact distinguishability and distinguishability at a prescribed accuracy need
not have the same dimension. An additive relation between observation
exponents can disappear under an arbitrarily small perturbation, creating
a new exact direction whose statistical amplitude tends to zero. Our
object is to compute both the exact information law and its resolved
finite-state counterpart through such a collision.

The likelihood accumulated by an experiment and the information needed
for its remaining decisions are different objects. A posterior family can
have many directions that no remaining experiment can test. Conversely,
a sparse one-step detector can generate a much larger collection of
future tests by multiplication. The appropriate question is therefore not
only the rank of the detector, but the rank of the pairing between an
\emph{attainable product tangent} and the \emph{future observation algebra}.

Here this pairing can be computed for a whole class of nonsaturated
experiments. Let
\[
 A=\{0=a_0<a_1<\cdots<a_{r-1}=D\},\qquad r\ge2,
 \quad W_A=\spanop\{t^a:a\in A\},
\]
and let $mA$ denote the set of sums of exactly $m$ members of $A$, with
$0A=\{0\}$. Thus zero-padding is allowed and $mA\subset(m+1)A$.
Write $h_A(m)=|mA|$. The parameter lies in a fixed compact interval
$I=[l,u]$ with $0\le l<u<\infty$, and its prior $\mu$ has full support.
The detector and comparator are specified in Section~\ref{sec:experiment}.
No function of the unknown parameter is supplied to the controller.

\begin{theorem}[Sparse finite-horizon information law]\label{thm:main}
Suppose the positive finite detector spans $W_A$ and every lookup entry
ranges over $[\eta,1-\eta]$, where $0<\eta<1/2$. For $n,m\ge0$ put
\begin{equation}\label{eq:main-dimension}
 d_A(n,m)=\min\{n(r-1),h_A(m)-1\}.
\end{equation}
Two attainable posteriors after $n$ reports have the same law under every
common future-only continuation of length $m$ if and only if their moments
at the exponents in $mA$ agree. The prediction map on all-failure command
histories has generic rank $d_A(n,m)$, and it has a smooth attainable patch
of this dimension with a continuous section to histories. The smallest
Euclidean dimension of a sufficient encoding continuous on each
known-report history stratum is exactly $d_A(n,m)$.

For a total budget $N$, these dimensions can be attained simultaneously
by a causal exact-real state. Its maximal size is
\begin{equation}\label{eq:peak}
 D_A(N)=\max_{0\le n\le N}d_A(n,N-n).
\end{equation}
Only prior moments with exponents in $NA$ are needed for its updates.
\end{theorem}

Generic rank here means rank outside a proper real algebraic exceptional
set in the finite command coordinates. It is not a claim of uniform
conditioning. The global upper encoding is an encoding of histories;
when it stores normalized factors it need not identify different
factorizations of the same posterior. We do not assert that every singular
quotient has a global coordinate chart of its local dimension.

For $A=\{0,1,\ldots,q\}$, \eqref{eq:main-dimension} becomes
$q\min(n,m)$. For $A=\{0,2,5\}$ it instead gives $d_A(3,2)=5$.
More generally, for $A=\{0,a,b\}\subset\N$ and
$B=b/\gcd(a,b)$,
\begin{equation}\label{eq:three-intro}
 h_A(m)=\binom{m+2}{2}-\one_{\{m\ge B\}}\binom{m-B+2}{2}.
\end{equation}
For incommensurable $a,b>0$, the same three-cell detector space has
$h_A(m)=\binom{m+2}{2}$ at every $m$. Thus neither one-step rank nor
maximum polynomial degree describes all the possible memory profiles.

The central step is a prior-independent transversality calculation, not
an assumption that the relevant projection has maximal rank. At factors
$1+c_i t^D$, with the $c_i$ distinct and small, the unnormalized tangent
is the monomial span with exponent set
\[
 B_{A,n}=\{jD:0\le j\le n\}
 \ \cup\!\!\bigcup_{a\in A\setminus\{0,D\}}
                  \{a+jD:0\le j<n\}.
\]
It has $n(r-1)+1$ members. Every square mixed-moment minor between this
set and $mA$ is strictly positive. Normalization removes exactly one rank.
Section~\ref{sec:rank} supplies the full argument, including attainment
inside one fixed comparator cube and the topological lower bound.


\subsection{The uniform law at an additive collision}
The exact formula is the starting point, rather than the end, of the
finite-resolution analysis. Let the exponents depend affinely on a known
calibration parameter, $A_\theta=\{a_i+\theta b_i\}$. Group the
formal $m$-fold sum exponents by their limiting value at zero. A cluster
of multiplicity $s$ contributes orders $0,1,\ldots,s-1$. After deleting
the constant coordinate, list all orders as
$\nu_1\le\cdots\le\nu_d$, with $d$ the nearby future dimension.
Theorem~\ref{thm:confluent-law} proves, whenever $d\le n(r-1)$,
\begin{equation}\label{eq:intro-confluent}
 R_{M;n,m}^\theta\asymp
 \max_{1\le\ell\le d}
       \theta^{2(\nu_1+\cdots+\nu_\ell)/\ell}M^{-2/\ell}
\end{equation}
for both unconditional and minimax checkpoint prediction, with constants
uniform through the collision. The calibration, prior and finite horizon
are fixed. Distinct positive collision parameters use the same command
and query labels. The observation of a logarithmic jet is not postulated:
Newton interpolation identifies its actual coefficient in the physical
query vector. Strict confluent positivity proves that the desingularized
attainable patch survives at zero, including for priors without a density.

The family $A_\theta=\{0,1,2+\theta\}$ supplies a complete sequential
realization of this law. In the positive three-cell model
\eqref{eq:near-collision-cells}, with total horizon five, its exact peak
is four at zero and five for every $0<\theta\le1/2$. Nevertheless the
new direction has width $\theta$, not unit width. Theorem~\ref{thm:uniform-streaming}
proves
\begin{equation}\label{eq:intro-uniform}
 \overline{\mathcal R}_{M,5}^{\theta,\mathrm{stream}}
   \asymp \max\{M^{-1/2},\theta^{2/5}M^{-2/5}\}
\end{equation}
uniformly for every integer $M\ge1$ and $\theta\in[0,1/2]$.
The same upper filter works on every admitted history and at every
checkpoint. The transition from a four-dimensional to a five-dimensional
rate occurs at order $M=\theta^{-4}$, or distortion $\theta^2$.
This determines, rather than leaves inside an unspecified constant, the
degeneration of the high-resolution law.

A singular coordinate system would not suffice to prove
\eqref{eq:intro-uniform}: a transition with norm $1/\theta$ could
amplify an earlier quantization error and destroy uniformity. The proof
uses normalized factors before the peak, four physical moments and one
unscaled moment difference at the peak, and two moments afterward. Its
shrinking update \eqref{eq:uniform-update} contains no division by
$\theta$. Reachable representatives therefore give a genuine index-only
filter with uniformly accumulated errors.

Finally, after the same exploration and query, offer a ticket on the
same future event at an independent uniform price. A binary accept-or-reject
action then has optimal value gap exactly half the squared-prediction
regret. More substantially, Theorem~\ref{thm:resolved-value} compares
one common payoff before and after the deterministic history statistic
$T_\theta$ discards the weak coordinate. The full loss is comparable
to $\theta^2$ uniformly, and a streaming controller with
$M\ge K\theta^{-4}$ beats even the uncompressed erased statistic by
at least a fixed positive multiple of $\theta^2$. This is not an
assertion that every four-dimensional encoding is inadequate, nor an
optimal-exploration theorem. It gives the resolved sparse direction a
quantitative value in one precisely specified sequential decision.

\subsection{Fixed-calibration consequences and mechanical realization}
An exact-real causal formula alone does not give a finite-state controller.
Section~\ref{sec:streaming} constructs transition maps that store only one
of $M$ symbols after \emph{every} report. For fixed $N\ge2$, its prediction
regret is at most $C_NM^{-2/D_A(N)}$ at any checkpoint. At a checkpoint
maximizing \eqref{eq:peak}, one fixed randomized exploration experiment
has regret at least $c_NM^{-2/D_A(N)}$ for every such finite-state filter,
including the declared independent randomization. These bounds match the
exponent, not the leading distortion constant. The query is announced
only after the checkpoint state has been formed, and old commands or
command-generation seeds cannot be retained outside that state.

The lower bound concerns a prescribed exploration-and-prediction task;
it is not a universal lower bound on all optimally chosen control tasks.
For bounded additive observable rewards we separately prove an exact-model
$C'_NM^{-1/D_A(N)}$ optimal-control loss bound. We keep that weaker
power distinct from the squared-prediction power. For an all-failure probe
whose factors are bounded by $v<1$, a one-state zero forecast has regret
at most $v^{2m}$. That elementary finite-resolution bound remains visible
alongside the new streaming rate.

Section~\ref{sec:common-risk} studies the collision bit using one common
prediction target and one common total-risk baseline. The optimum for
\emph{every} memory budget is given by a finite partition formula. Four
states realize the full strict gain. For a sufficiently small positive
amplitude, three states have exactly the same optimal risk as the erased
experiment, even though the exact detector rank is already larger. This
identifies an operational resolution effect rather than equating rank with
useful accuracy.

\subsection{Relation to earlier work}
Action-conditional predictions as representations of state are classical
\cite{LittmanSutton}. Approximate recursively updated information states
and their dynamic-programming loss bounds are treated in
\cite{Subramanian}; finite-window policies under filter stability are
studied in \cite{KaraYuksel}. Squared-error quantization supplies the
covering exponent \cite{GrafLuschgy}. Total positivity and generalized
Vandermonde matrices are classical \cite{Pinkus}. The identification of
iterated sumsets with Hilbert functions is also established
\cite{EliahouMazumdar,Elias}; it is not a novelty claim here.

The result proved here concerns their specific junction: the attainable
normalized-product tangent of a sparse calibrated experiment pairs at
maximal rank with its future test space for every full-support prior.
This yields the explicit law \eqref{eq:main-dimension}, its causal
realization, and the finite-state consequences with the same attainable
geometry. None of the cited general representation, sumset, positivity or
quantization principles alone is used as a substitute for this calculation.
The complete v6 predecessor is retained \cite{A1v6}. The full-space
predecessor and its mechanical derivations are retained
as a complete source companion \cite{A1v5}. The sparse three-cell example
that motivated the present generalization was supplied in the source-pinned
v5 referee report \cite{A1review}; its authorship is not reassigned here.

The additional quantitative input is the uniform attainable pairing with
confluent exponent jets and its transfer to actual sequential state
updates. The divided-difference identities and rectangular covering
calculation are classical; they do not by themselves establish attainable
surjectivity at the singular calibration, or prevent error amplification
in a causal implementation. The near-collision example and its preliminary
base-model bound originate in \cite{A1v6note}. Equations
\eqref{eq:intro-confluent}--\eqref{eq:intro-uniform} and the same-task
weak-direction value are the assertions whose proofs occupy
Sections~\ref{sec:confluence}, \ref{sec:uniform} and
\ref{sec:resolved-value}.
